% ==============================================================================
% PRESENTACIÓN: Equidad en Aprendizaje Automático — Bank Marketing
% Compilar con: pdflatex presentacion.tex (2 veces para TOC)
% ==============================================================================
\documentclass[aspectratio=169, 11pt]{beamer}

% --- TEMA Y COLORES ---
\usetheme{default}
\useinnertheme{rounded}
\useoutertheme{miniframes}

% Paleta: grayscale + acento azul oscuro
\definecolor{azuloscuro}{HTML}{1A3A5C}
\definecolor{azulmedio}{HTML}{2E6DA4}
\definecolor{azulclaro}{HTML}{87B7DB}
\definecolor{gris90}{HTML}{1A1A1A}
\definecolor{gris70}{HTML}{4D4D4D}
\definecolor{gris50}{HTML}{808080}
\definecolor{gris30}{HTML}{B3B3B3}
\definecolor{gris10}{HTML}{E6E6E6}
\definecolor{fondoslide}{HTML}{F5F5F5}
\definecolor{cardcolor}{HTML}{FFFFFF}

\setbeamercolor{palette primary}{bg=azuloscuro, fg=white}
\setbeamercolor{palette secondary}{bg=azulmedio, fg=white}
\setbeamercolor{palette tertiary}{bg=gris70, fg=white}
\setbeamercolor{structure}{fg=azuloscuro}
\setbeamercolor{title}{fg=white, bg=azuloscuro}
\setbeamercolor{frametitle}{fg=azuloscuro, bg=gris10}
\setbeamercolor{background canvas}{bg=fondoslide}
\setbeamercolor{block title}{bg=azuloscuro, fg=white}
\setbeamercolor{block body}{bg=cardcolor, fg=gris90}
\setbeamercolor{block title alerted}{bg=azulmedio, fg=white}
\setbeamercolor{block body alerted}{bg=cardcolor, fg=gris70}
\setbeamercolor{itemize item}{fg=azuloscuro}
\setbeamercolor{itemize subitem}{fg=azulmedio}
\setbeamercolor{normal text}{fg=gris90}
\setbeamercolor{footline}{fg=gris50}

% Puntas redondeadas
\setbeamertemplate{blocks}[rounded][shadow=false]
\setbeamertemplate{itemize items}[circle]
\setbeamertemplate{navigation symbols}{}
\setbeamertemplate{footline}{%
  \hfill\insertframenumber/\inserttotalframenumber\hspace*{6pt}\vspace{4pt}%
}

% --- TIPOGRAFÍA ---
\usepackage[T1]{fontenc}
\usepackage[utf8]{inputenc}
\usepackage[spanish]{babel}
\usepackage{lmodern}
\usepackage{microtype}

% --- PAQUETES GRÁFICOS ---
\usepackage{graphicx}
\usepackage{tikz}
\usetikzlibrary{shapes.geometric, arrows.meta, positioning, calc, shadows.blur}
\usepackage{booktabs}
\usepackage{colortbl}
\usepackage{multirow}
\usepackage{fontawesome5}

% --- METADATOS ---
\title[Equidad en ML]{Equidad en Aprendizaje Automático}
\subtitle{Trabajo Práctico — Bank Marketing Dataset}
\author{Grupo TP}
\date{1er Cuatrimestre 2026}
\institute{Equidad en Aprendizaje Automático}

% ==============================================================================
\begin{document}

% ---- PORTADA ----
{
\setbeamercolor{background canvas}{bg=azuloscuro}
\begin{frame}[plain]
\vfill
\begin{center}
{\color{white}\Huge\bfseries Equidad en\\Aprendizaje Automático}\\[12pt]
{\color{azulclaro}\large Trabajo Práctico — Bank Marketing Dataset}\\[24pt]
{\color{gris30}\normalsize 1er Cuatrimestre 2026}\\[6pt]
\end{center}
\vfill
\end{frame}
}

% ---- ÍNDICE ----
\begin{frame}{Hoja de ruta}
\begin{columns}[T]
\begin{column}{0.48\textwidth}
\begin{block}{\faSearch\quad Exploración}
\begin{itemize}
\item[\textbf{Ej1}] Conocer el dataset
\item[\textbf{Ej2}] Modelo baseline
\end{itemize}
\end{block}
\vspace{4pt}
\begin{block}{\faBalanceScale\quad Auditoría}
\begin{itemize}
\item[\textbf{Ej3}] Criterios de fairness
\end{itemize}
\end{block}
\end{column}
\begin{column}{0.48\textwidth}
\begin{block}{\faWrench\quad Mitigación}
\begin{itemize}
\item[\textbf{Ej4}] Técnicas de corrección
\end{itemize}
\end{block}
\vspace{4pt}
\begin{block}{\faChartBar\quad Evaluación}
\begin{itemize}
\item[\textbf{Ej5}] Comparación y reflexión
\end{itemize}
\end{block}
\end{column}
\end{columns}
\end{frame}

% ==============================================================================
% EJERCICIO 1
% ==============================================================================
{
\setbeamercolor{background canvas}{bg=azuloscuro}
\begin{frame}[plain]
\vfill
\begin{center}
{\color{white}\LARGE\bfseries Ejercicio 1}\\[8pt]
{\color{azulclaro}\large Conocer el dataset}
\end{center}
\vfill
\end{frame}
}

\begin{frame}{El dataset: Bank Marketing}
\begin{columns}[T]
\begin{column}{0.55\textwidth}
\begin{block}{Ficha técnica}
\begin{itemize}
\item \textbf{Fuente:} UCI ML Repository
\item \textbf{Registros:} 41.188 contactos telefónicos
\item \textbf{Features:} 20 + target (\texttt{y})
\item \textbf{Target:} ¿Suscripción a plazo fijo?
\item \textbf{Periodo:} 2008--2013 (Portugal)
\end{itemize}
\end{block}
\vspace{4pt}
\begin{alertblock}{Variables protegidas}
\begin{itemize}
\item \textbf{Estado civil} (\texttt{marital})
\item \textbf{Edad} (\texttt{age})
\item \textbf{Género} → sin variable directa, usamos \texttt{job} como proxy
\end{itemize}
\end{alertblock}
\end{column}
\begin{column}{0.42\textwidth}
\begin{block}{Autores}
Moro, Cortez \& Rita (2014). Datos de un banco portugués bajo confidencialidad.
\end{block}
\vspace{4pt}
\begin{block}{\faExclamationTriangle\quad Data leakage}
La variable \texttt{duration} (duración de la llamada) se conoce \emph{después} del hecho. Los propios autores advierten excluirla para modelos predictivos reales.\\[4pt]
\textbf{→ La excluimos desde el inicio.}
\end{block}
\end{column}
\end{columns}
\end{frame}

\begin{frame}{Sesgos identificados en los datos}
\begin{columns}[T]
\begin{column}{0.48\textwidth}
\begin{block}{\faChartPie\quad Desbalanceo de clases}
\begin{itemize}
\item \textasciitilde89\% de los contactos → \texttt{no}
\item Solo \textasciitilde11\% → \texttt{yes}
\item El modelo puede ``aprender'' a decir no siempre
\end{itemize}
\end{block}
\vspace{4pt}
\begin{block}{\faUsers\quad Estado civil}
\begin{itemize}
\item \texttt{married}: \textasciitilde60\% (sobrerep.)
\item \texttt{single}: \textasciitilde28\%
\item \texttt{divorced}: \textasciitilde11\% (subrep.)
\item \texttt{single} muestra señal positiva
\end{itemize}
\end{block}
\end{column}
\begin{column}{0.48\textwidth}
\begin{block}{\faBirthdayCake\quad Edad}
\begin{itemize}
\item Jóvenes y jubilados subrepresentados
\item Jubilados tienen la tasa de \texttt{yes} \textbf{más alta} (\textasciitilde27\%)
\item Pero pocos datos → modelo los ignora
\item \textbf{Paradoja de la tasa base}
\end{itemize}
\end{block}
\vspace{4pt}
\begin{block}{\faBriefcase\quad Job como proxy de género}
\begin{itemize}
\item \texttt{housemaid}, \texttt{admin.} → históricamente femeninos
\item Sin variable género directa
\end{itemize}
\end{block}
\end{column}
\end{columns}
\end{frame}

% ==============================================================================
% EJERCICIO 2
% ==============================================================================
{
\setbeamercolor{background canvas}{bg=azuloscuro}
\begin{frame}[plain]
\vfill
\begin{center}
{\color{white}\LARGE\bfseries Ejercicio 2}\\[8pt]
{\color{azulclaro}\large Modelo baseline}
\end{center}
\vfill
\end{frame}
}

\begin{frame}{Modelo: Random Forest}
\begin{columns}[T]
\begin{column}{0.48\textwidth}
\begin{block}{Pipeline}
\begin{enumerate}
\item Eliminar \texttt{duration} (data leakage)
\item Filtrar \texttt{marital = unknown} (80 reg.)
\item Codificar categóricas (LabelEncoder)
\item Split estratificado 80/20
\item Random Forest (100 estimadores)
\end{enumerate}
\end{block}
\end{column}
\begin{column}{0.48\textwidth}
\begin{block}{Resultados globales}
\begin{center}
\begin{tabular}{l r}
\toprule
\textbf{Métrica} & \textbf{Valor} \\
\midrule
\rowcolor{gris10} Accuracy & 90\% \\
Precision (yes) & 60\% \\
\rowcolor{gris10} \textbf{Recall (yes)} & \textbf{30\%} \\
F1 (yes) & 40\% \\
\bottomrule
\end{tabular}
\end{center}
\vspace{4pt}
{\small\color{gris50} Accuracy engañoso por desbalanceo.\\
Recall bajo = muchos FN}
\end{block}
\end{column}
\end{columns}
\end{frame}

\begin{frame}{¿Cuál error es peor?}
\begin{center}
\begin{tikzpicture}[
  card/.style={rounded corners=8pt, draw=gris30, fill=cardcolor, 
    minimum width=6.5cm, minimum height=2.2cm, align=center,
    blur shadow={shadow blur steps=5, shadow xshift=1pt, shadow yshift=-1pt}},
]
\node[card] (fn) at (0,0) {
  \textbf{\large\color{azuloscuro} Falso Negativo (FN)}\\[4pt]
  {\small Cliente que se suscribiría}\\
  {\small pero el modelo dice ``no'' → \textbf{NO lo llamamos}}
};
\node[card] (fp) at (8.5,0) {
  \textbf{\large\color{gris50} Falso Positivo (FP)}\\[4pt]
  {\small Cliente que no se suscribiría}\\
  {\small pero el modelo dice ``sí'' → una llamada extra}
};
\node[below=0.4cm of fn, text=azuloscuro, font=\bfseries] {PEOR: pérdida de negocio};
\node[below=0.4cm of fp, text=gris50] {Costo marginal bajo};
\end{tikzpicture}
\end{center}
\vspace{6pt}
\begin{block}{Justificación}
\begin{itemize}
\item El \textbf{FN} pierde directamente un depósito y un cliente potencial
\item El \textbf{FP} cuesta solo una llamada telefónica adicional
\item Objetivo del banco: \emph{incrementar suscripciones} → priorizar Recall
\end{itemize}
\end{block}
\end{frame}

% ==============================================================================
% EJERCICIO 3
% ==============================================================================
{
\setbeamercolor{background canvas}{bg=azuloscuro}
\begin{frame}[plain]
\vfill
\begin{center}
{\color{white}\LARGE\bfseries Ejercicio 3}\\[8pt]
{\color{azulclaro}\large Auditoría de fairness}
\end{center}
\vfill
\end{frame}
}

\begin{frame}{Criterios de fairness — en contexto}
\begin{columns}[T]
\begin{column}{0.48\textwidth}
\begin{block}{Statistical Parity}
Igual \textbf{proporción de predicciones positivas} entre grupos, independientemente del resultado real.
\end{block}
\vspace{4pt}
\begin{block}{Equal Opportunity}
Igual \textbf{Recall (TPR)} entre grupos. Detectar suscriptores igual de bien para todos.
\end{block}
\end{column}
\begin{column}{0.48\textwidth}
\begin{block}{Equalized Odds}
Igual \textbf{TPR + FPR} entre grupos. Más exigente: equidad tanto en aciertos como en falsas alarmas.
\end{block}
\vspace{4pt}
\begin{block}{Predictive Parity}
Igual \textbf{Precision} entre grupos. La confianza de cada predicción positiva es la misma.
\end{block}
\end{column}
\end{columns}
\vspace{8pt}
\begin{alertblock}{\faExclamationCircle\quad Resultado teórico clave (Chouldechova, 2017)}
Cuando las tasas base difieren entre grupos, es \textbf{imposible} satisfacer los tres criterios simultáneamente. \textbf{Debemos elegir uno.}
\end{alertblock}
\end{frame}

\begin{frame}{Fairness del modelo base — proxy de género (\texttt{job})}
\begin{center}
\begin{tabular}{l c c c}
\toprule
\textbf{Métrica} & \textbf{Hist. femenino} & \textbf{Hist. masc./otro} & \textbf{Disparidad} \\
\midrule
\rowcolor{gris10} Statistical Parity & 0.055 & 0.053 & 0.002 \\
\textbf{Equal Opportunity (TPR)} & \textbf{0.315} & \textbf{0.296} & \textbf{0.019} \\
\rowcolor{gris10} FPR (Eq. Odds) & 0.026 & 0.025 & 0.001 \\
Predictive Parity (Prec.) & 0.648 & 0.583 & 0.065 \\
\bottomrule
\end{tabular}
\end{center}
\vspace{6pt}
\begin{columns}[T]
\begin{column}{0.48\textwidth}
\begin{block}{Umbral elegido: 0.10}
Disparidades $< 0.10$ → modelo es \textbf{fair}.\\
Todas las disparidades están por debajo → el modelo base ya es equitativo para este proxy.
\end{block}
\end{column}
\begin{column}{0.48\textwidth}
\begin{alertblock}{Criterio elegido: Equal Opportunity}
Coherente con Ej2: si FN es el error más costoso, queremos que el \textbf{Recall} sea igual entre grupos → no perder suscriptores de ningún grupo.
\end{alertblock}
\end{column}
\end{columns}
\end{frame}

\begin{frame}{Fairness del modelo base — estado civil (\texttt{marital})}
\begin{center}
\begin{tabular}{l c c c c}
\toprule
\textbf{Métrica} & \textbf{Married} & \textbf{Single} & \textbf{Divorced} & \textbf{Max Disp.} \\
\midrule
\rowcolor{gris10} TPR (Equal Opp.) & 0.274 & 0.348 & 0.294 & 0.075 \\
FPR (Eq. Odds) & --- & --- & --- & 0.020 \\
\bottomrule
\end{tabular}
\end{center}
\vspace{6pt}
\begin{block}{Observación}
\begin{itemize}
\item Disparidad TPR = 0.075 → \textbf{por debajo del umbral 0.10}, pero más alta que en \texttt{job}
\item \texttt{single} tiene el TPR más alto (más jóvenes, señal positiva)
\item \texttt{married} tiene el TPR más bajo (grupo mayoritario, comportamiento conservador)
\end{itemize}
\end{block}
\end{frame}

% ==============================================================================
% EJERCICIO 4
% ==============================================================================
{
\setbeamercolor{background canvas}{bg=azuloscuro}
\begin{frame}[plain]
\vfill
\begin{center}
{\color{white}\LARGE\bfseries Ejercicio 4}\\[8pt]
{\color{azulclaro}\large Mitigación de sesgos}
\end{center}
\vfill
\end{frame}
}

\begin{frame}{Técnicas de mitigación aplicadas}
\begin{columns}[T]
\begin{column}{0.48\textwidth}
\begin{block}{\faBalanceScale\quad Reweighting (Pre-processing)}
\begin{itemize}
\item Asigna \textbf{pesos} a las muestras de entrenamiento
\item Compensa sub-representación de grupos
\item Actúa \emph{antes} de entrenar el modelo
\item El modelo aprende ``dando más importancia'' a los grupos rezagados
\end{itemize}
\end{block}
\end{column}
\begin{column}{0.48\textwidth}
\begin{block}{\faSlidersH\quad Threshold Adjustment (Post-processing)}
\begin{itemize}
\item Ajusta el \textbf{umbral de decisión por grupo}
\item Iguala directamente el TPR (Equal Opportunity)
\item Actúa \emph{después} de entrenar el modelo
\item Más directo para el criterio elegido
\end{itemize}
\end{block}
\end{column}
\end{columns}
\vspace{8pt}
\begin{center}
\begin{tikzpicture}[
  box/.style={rounded corners=6pt, minimum width=3cm, minimum height=1cm, 
    align=center, font=\small\bfseries},
  arr/.style={-{Stealth[length=6pt]}, thick, color=azulmedio}
]
\node[box, fill=gris10, draw=gris30] (data) at (0,0) {Datos};
\node[box, fill=azulclaro!30, draw=azulmedio] (rw) at (3.5,0) {Reweighting};
\node[box, fill=gris10, draw=gris30] (model) at (7,0) {Random Forest};
\node[box, fill=azulclaro!30, draw=azulmedio] (ta) at (10.5,0) {Threshold Adj.};
\node[box, fill=gris10, draw=gris30] (pred) at (14,0) {Predicción};
\draw[arr] (data) -- (rw);
\draw[arr] (rw) -- (model);
\draw[arr] (model) -- (ta);
\draw[arr] (ta) -- (pred);
\end{tikzpicture}
\end{center}
\end{frame}

% ==============================================================================
% EJERCICIO 5
% ==============================================================================
{
\setbeamercolor{background canvas}{bg=azuloscuro}
\begin{frame}[plain]
\vfill
\begin{center}
{\color{white}\LARGE\bfseries Ejercicio 5}\\[8pt]
{\color{azulclaro}\large Comparación y reflexión}
\end{center}
\vfill
\end{frame}
}

\begin{frame}{Comparación de modelos — Performance global}
\begin{center}
\begin{tabular}{l c c c c}
\toprule
\textbf{Modelo} & \textbf{Accuracy} & \textbf{Precision} & \textbf{Recall} & \textbf{F1} \\
\midrule
\rowcolor{gris10} Original & 0.899 & 0.603 & 0.302 & 0.403 \\
Reweighting & 0.897 & 0.593 & 0.283 & 0.383 \\
\rowcolor{gris10} Threshold Adj. & 0.897 & 0.583 & \textbf{0.316} & \textbf{0.410} \\
\bottomrule
\end{tabular}
\end{center}
\vspace{6pt}
\begin{columns}[T]
\begin{column}{0.48\textwidth}
\begin{block}{Observación}
\begin{itemize}
\item Accuracy casi constante (\textasciitilde90\%)
\item Threshold Adj. mejora Recall (+1.4pp) y F1
\item Reweighting pierde Recall → no funciona tan bien aquí
\end{itemize}
\end{block}
\end{column}
\begin{column}{0.48\textwidth}
\begin{alertblock}{Trade-off}
Pequeña pérdida en Precision a cambio de mayor Recall.\\[4pt]
\textbf{Para el banco:} algunas llamadas extra (FP) pero más suscriptores captados (menos FN).
\end{alertblock}
\end{column}
\end{columns}
\end{frame}

\begin{frame}{Comparación de fairness — Equal Opportunity}
\begin{center}
\textbf{Proxy de género (\texttt{job})}\\[4pt]
\begin{tabular}{l c c c}
\toprule
\textbf{Modelo} & \textbf{TPR fem.} & \textbf{TPR masc.} & \textbf{Disp. TPR} \\
\midrule
\rowcolor{gris10} Original & 0.315 & 0.296 & 0.019 \\
Reweighting & 0.315 & 0.293 & 0.022 \\
\rowcolor{gris10} Threshold Adj. & 0.348 & 0.302 & 0.046 \\
\bottomrule
\end{tabular}
\end{center}
\vspace{8pt}
\begin{center}
\textbf{Estado civil (\texttt{marital})}\\[4pt]
\begin{tabular}{l c c c c}
\toprule
\textbf{Modelo} & \textbf{Married} & \textbf{Single} & \textbf{Divorced} & \textbf{Max Disp.} \\
\midrule
\rowcolor{gris10} Original & 0.274 & 0.348 & 0.294 & 0.075 \\
Reweighting & 0.244 & 0.345 & 0.271 & 0.101 \\
\rowcolor{gris10} Threshold Adj. & 0.344 & 0.345 & 0.365 & \textbf{0.020} \\
\bottomrule
\end{tabular}
\end{center}
\vspace{4pt}
{\small Threshold Adj. reduce la disparidad de \textbf{0.075 → 0.020} para \texttt{marital}. Gran mejora.}
\end{frame}

\begin{frame}{Reflexión: impacto en el mundo real}
\begin{columns}[T]
\begin{column}{0.48\textwidth}
\begin{block}{\faUniversity\quad Acceso equitativo}
Sin mitigación, ciertos grupos quedan excluidos de ofertas financieras. El modelo perpetúa sesgos históricos.
\end{block}
\vspace{4pt}
\begin{block}{\faGavel\quad Regulación}
La discriminación por género, edad o estado civil en servicios financieros es ilegal en la mayoría de jurisdicciones.
\end{block}
\end{column}
\begin{column}{0.48\textwidth}
\begin{block}{\faCogs\quad Trade-off consciente}
No existe modelo justo Y perfecto. La elección del criterio de fairness es una decisión de \emph{política}, no solo técnica.
\end{block}
\vspace{4pt}
\begin{block}{\faExclamationTriangle\quad Limitaciones}
\texttt{job} como proxy de género es una aproximación. No todas las personas en \texttt{admin.} son mujeres. Datos demográficos reales mejorarían la auditoría.
\end{block}
\end{column}
\end{columns}
\end{frame}

% ---- CIERRE ----
{
\setbeamercolor{background canvas}{bg=azuloscuro}
\begin{frame}[plain]
\vfill
\begin{center}
{\color{white}\LARGE\bfseries Conclusión}\\[12pt]
\begin{tikzpicture}[
  step/.style={rounded corners=6pt, minimum width=2.2cm, minimum height=0.9cm,
    align=center, font=\small\bfseries, fill=azulmedio, text=white},
  arr/.style={-{Stealth[length=5pt]}, thick, color=azulclaro}
]
\node[step] (e1) at (0,0) {Ej1\\Datos};
\node[step] (e2) at (3,0) {Ej2\\Baseline};
\node[step] (e3) at (6,0) {Ej3\\Auditoría};
\node[step] (e4) at (9,0) {Ej4\\Mitigación};
\node[step] (e5) at (12,0) {Ej5\\Evaluación};
\draw[arr] (e1) -- (e2);
\draw[arr] (e2) -- (e3);
\draw[arr] (e3) -- (e4);
\draw[arr] (e4) -- (e5);
\end{tikzpicture}
\vspace{16pt}

{\color{azulclaro}\large Pipeline completo de auditoría y mitigación de sesgos\\[6pt]
aplicable a cualquier sistema de ML en producción.}
\end{center}
\vfill
\end{frame}
}

\end{document}
